\documentclass[12pt,a4paper]{article}

\usepackage[spanish,es-nodecimaldot]{babel}
\usepackage[utf8]{inputenc}
\usepackage[T1]{fontenc}
\usepackage{lmodern}
\usepackage{amsmath,amssymb}
\usepackage{graphicx}
\usepackage{booktabs}
\usepackage{hyperref}
\usepackage{enumitem}
\usepackage{geometry}
\usepackage{listings}
\usepackage{xcolor}

\geometry{margin=2.5cm}

\lstset{
  basicstyle=\footnotesize\ttfamily,
  breaklines=true,
  frame=single,
}

\begin{document}
\section{Deep research para defender TFG\_CYBER\_AI}

\subsection{Qué he decidido investigar}

He dejado fuera, como pediste, la investigación interna de \texttt{sb3\_contrib} y de QR-DQN. Tras
inspeccionar el repositorio en su estado actual sobre \texttt{main}, lo más importante para tu
defensa no es volver a explicar el algoritmo, sino dominar el \textbf{contrato de datos} del
proyecto, el \textbf{adapter de CICIDS2017}, la \textbf{formalización del problema como entorno RL},
el \textbf{pipeline de entrenamiento y reproducibilidad}, la \textbf{validación} y la \textbf{Fase~2
de inferencia sobre tráfico real}. Esos son, de hecho, los entry points y los invariantes que el
propio repo declara como ``source of truth'': \verb|src/canonical_schema.py|,
\verb|src/load_cicids2017.py|, \verb|src/rl_defender_env.py|, \verb|src/train_rl_defender.py|,
\verb|src/validate_checks.py|, \verb|src/validate_leave_one_csv_out.py| y
\verb|scripts/predict_real_traffic_v2.py|.

Mi lectura está orientada a que puedas responder no solo \textit{qué has usado}, sino
\textbf{por qué existe cada pieza}, \textbf{qué hipótesis codifica}, \textbf{qué garantiza},
\textbf{qué no garantiza} y \textbf{qué limitaciones honestas tiene hoy el sistema}. El foco
correcto, por tanto, no es ``tengo un modelo que clasifica bien'', sino ``tengo un pipeline
reproducible y defendible desde el dataset hasta el artefacto de Fase~2''. Eso también coincide con
las notas de defensa ya presentes en el repo.

\subsection{Esquema canónico y contrato de datos}

La decisión arquitectónica más importante del proyecto no es el algoritmo, sino el \textbf{espacio
de observación estable}. El dataset oficial CICIDS2017, publicado por el Canadian Institute for
Cybersecurity, ofrece PCAPs y CSVs con \textbf{más de 80 features de flujo} extraídas con
CICFlowMeter, a lo largo de cinco días de tráfico normal y malicioso. Tu repo no usa ``todo lo que
haya'', sino que selecciona \textbf{76 features canónicas flow-based}, deliberadamente numéricas,
estables y extraíbles también desde tráfico real; además excluye IPs, timestamps, Flow IDs y puertos
específicos por riesgo de leakage o de representar atajos espurios hacia la etiqueta.

Ese diseño produce un vector final de \textbf{152 dimensiones}: las 76 features canónicas y otras 76
posiciones de máscara de missingness. La idea de añadir indicadores binarios de ausencia es sólida
desde el punto de vista de ingeniería de datos: en aprendizaje automático, los \textbf{missing
indicators} se usan precisamente para no confundir ``valor imputado'' con ``valor observado'', de
modo que el modelo pueda explotar el patrón de ausencia como señal adicional. En tu código, la
semántica declarada es \texttt{1 = present/valid} y \texttt{0 = missing/imputed}.

Aquí hay una sutileza muy importante para tu defensa: \textbf{mi lectura del código sugiere que, en
CICIDS2017, la máscara no conserva todos los NaN originales fila a fila}. El motivo es que
\verb|load_cicids2017.py| reemplaza \texttt{inf} por \texttt{NaN} y luego rellena las features con
\texttt{0} antes de llamar a \verb|map_to_canonical()|. Después, \verb|map_to_canonical()| marca la
máscara como 0 solo si todavía ve valores no finitos. Eso significa que, en la práctica del pipeline
actual, la máscara parece servir sobre todo para \textbf{ausencias estructurales} entre datasets o
extractores, no para preservar toda la missingness original de cada celda de CICIDS2017. Es una
observación muy defendible si te preguntan ``¿la máscara codifica missingness real del dataset o
compatibilidad de esquema?''; la respuesta correcta hoy es: \textbf{principalmente compatibilidad
de esquema y extractor, según el flujo actual de preprocesado}.

Otra sutileza todavía más fina: el \texttt{StandardScaler} se ajusta sobre el vector completo de 152
dimensiones, no solo sobre las 76 features reales. Eso implica que la máscara entra al modelo
\textbf{escalada}, no como bits crudos 0/1. Conceptualmente, esto transforma la señal
``presente/ausente'' en ``presente/ausente respecto al patrón estadístico de train''. Si en
entrenamiento una mask feature es casi siempre 1, tras escalar tenderá a 0; si en inferencia aparece
un 0 donde antes casi nunca lo había, entrará como una desviación clara. En otras palabras: aunque
sobre el papel la máscara sea binaria, \textbf{a la red le llega como feature estandarizada}, y eso
refuerza su papel como detector de mismatch estructural.

La defensa verbal más fuerte aquí es esta: \textbf{no has diseñado un modelo acoplado a un CSV
concreto}, sino un contrato de observación estable entre datasets históricos y tráfico real. Esa es,
probablemente, la decisión más ``de ingeniería de sistema'' de todo el TFG.

\subsection{Carga de CICIDS2017 y preprocesado}

El adapter de CICIDS2017 (\verb|src/load_cicids2017.py|) está más trabajado de lo que parece a
simple vista. Reconoce explícitamente los \textbf{8 CSVs oficiales} del dataset, lee por
\textit{chunks} para no romper RAM, localiza la columna de etiqueta de forma robusta, binariza
\texttt{BENIGN} frente a \texttt{ATTACK}, y elimina columnas tipo identificador como
\texttt{Flow~ID}, \texttt{Timestamp}, IPs y puertos. Esta eliminación no es cosmética: scikit-learn
recuerda que hay leakage cuando el modelo aprende con información que no estará disponible ---o no
debería influir causalmente--- en el momento real de predicción; los identificadores y proxies
temporales son candidatos clásicos a ese problema.

Tienes tres formas de particionar CICIDS2017, y deberías saber defender las tres. La primera es el
\textbf{random stratified split}, útil para iteración rápida y para cuantificar rendimiento offline
bajo distribución similar. La segunda es el \textbf{day/CSV split}, donde se entrena en unos
días/CSV y se evalúa en otros; esto es más realista porque el dataset oficial está organizado por
días con condiciones y ataques distintos. La tercera es el \textbf{leave-one-exact-CSV-out}, que
deja fuera exactamente uno de los CSV reales y entrena sobre los otros siete. Esa tercera vía es
metodológicamente la más fina porque reduce el riesgo de que el modelo ``vea'' patrones demasiado
parecidos entre train y test.

La lógica del day split se entiende mejor si recuerdas cómo está construido CICIDS2017: la página
oficial describe cinco días de captura, con \textbf{lunes de actividad normal} y \textbf{martes a
viernes con tráfico benigno y ataques}. Por eso tiene sentido que el repo use, por defecto,
lunes-martes-miércoles para train y jueves-viernes para test en el modo ``day'': no es un capricho,
sino un intento de medir generalización a patrones temporales distintos. Si te preguntan ``¿por qué
no basta con un split aleatorio?'', la respuesta fuerte es: \textbf{porque un split aleatorio puede
sobreestimar rendimiento cuando train y test comparten distribuciones demasiado parecidas; el split
por día/CSV es más exigente y más cercano al escenario de despliegue}.

También es correcta la forma en que el repo escala los datos: el \texttt{StandardScaler} se ajusta
sobre train y después se aplica a test, justo como recomiendan las guías de buenas prácticas de
scikit-learn para evitar contaminación entre subconjuntos. Además, la documentación oficial de
\texttt{StandardScaler} advierte que este transformador es \textbf{sensible a outliers}, un punto
crucial para entender por qué en Fase~2 aparecen percentiles y clipping en z-score.

\subsection{Entorno RL y qué problema está resolviendo realmente}

Tu entorno \texttt{RLDatasetDefenderEnv} es conceptualmente simple y, precisamente por eso, hay que
saber explicarlo muy bien. La observación es la muestra actual \(X[i]\). La acción es binaria:
\texttt{0 = PERMIT}, \texttt{1 = BLOCK}. La recompensa se calcula \textbf{solo} a partir de la
etiqueta real de esa muestra y de la acción elegida; después el entorno avanza a la siguiente
muestra. No hay estado latente de red, no hay evolución del atacante en respuesta a la acción del
agente, no hay dinámica compleja de sistema. Lo que hay es una secuencia de decisiones por flujo con
coste asimétrico.

Eso te lleva a una respuesta muy importante si el tribunal te pregunta ``¿esto es de verdad RL o es
clasificación disfrazada?''. La respuesta honesta es: \textbf{en la implementación actual, el
problema se parece mucho a un contextual bandit o a una clasificación secuencial con coste}, porque
en cada paso el agente recibe contexto (el flujo), elige acción y recibe recompensa basada en
contexto+acción. Esa es, de hecho, una definición estándar de contextual bandit en la literatura. La
parte defendible a favor de RL no está en decir que hoy ya modelas una red interactiva completa,
sino en decir que \textbf{la formulación por recompensa te permite incorporar directamente la
asimetría de costes de ciberseguridad y deja abierta la evolución futura hacia bloqueo activo y
decisiones realmente secuenciales}.

La recompensa actual del repo es \texttt{tp=1.5}, \texttt{fp=-1.5}, \texttt{fn=-5.0},
\texttt{omission=0.0}, y esto está alineado en entrenamiento, validación y entorno. Pero aquí hay
otra sutileza de código que conviene dominar: aunque el docstring del entorno todavía habla de un
posible \texttt{tn}, \textbf{la implementación ejecutada no usa una recompensa separada para true
negative}. Cuando el tráfico es benigno y la acción es \texttt{PERMIT}, el retorno es
\texttt{omission}, que por defecto vale 0.0. Es decir, el agente recibe una fuerte penalización por
dejar pasar ataques, otra penalización por bloquear benignos, una recompensa moderada por bloquear
ataques, y neutralidad por permitir benignos. En términos de diseño, eso prioriza no cometer falsos
negativos y no convierte el PERMIT correcto en una ``recompensa positiva'', sino en ausencia de
castigo. Debes poder explicarlo así, porque es exactamente la lógica codificada.

También debes conocer la mecánica de episodios. En entrenamiento, \verb|train_rl_defender.py| fija
\texttt{max\_steps\_ep = min(10\_000, len(X\_train))}, y cada \texttt{reset()} baraja los índices.
Eso significa que un episodio no equivale necesariamente a ``recorrer entero el dataset''; más bien,
el entrenamiento consume \textbf{segmentos aleatorizados} del conjunto de train. En evaluación, en
cambio, \texttt{shuffle=False} y el test se recorre en orden determinista. Y Check~A, además,
elimina por completo la dependencia del entorno porque compara \verb|model.predict(X_test[i])|
directamente contra \verb|y_test[i]|. Esto es importante para defender que las métricas no dependen
de una magia interna del \texttt{env.step()}.

Finalmente, hay dos piezas de \textit{framework plumbing} que debes saber justificar.
\texttt{DummyVecEnv} existe porque Stable-Baselines3 trabaja con entornos vectorizados y, según su
propia documentación, puede usarse incluso cuando solo quieres entrenar con un único entorno.
\texttt{Monitor}, por su parte, envuelve el entorno para registrar recompensas, duración y tiempo
de episodio. Eso no es ``decoración'': es parte de la trazabilidad experimental.

\subsection{Entrenamiento, artefactos y reproducibilidad}

El script de entrenamiento actual hace varias cosas bien desde el punto de vista de reproducibilidad.
Primero, crea un \texttt{RUN\_ID} con timestamp y lo usa para organizar salidas. Segundo, carga
CICIDS2017 \textbf{sin escalar} al principio para poder calcular percentiles sobre train y persistir
el \texttt{scaler} correctamente. Tercero, guarda no solo el \texttt{.zip} del modelo, sino también
\texttt{config.json}, \texttt{metrics.json}, \texttt{scaler.joblib} y
\texttt{train\_percentiles.npz}. Eso convierte cada experimento en un artefacto autocontenido y
reutilizable en Fase~2.

Hay además dos detalles finos que conviene saber. El primero es que el script calcula percentiles
\(p_{0.5}\) y \(p_{99.5}\) \textbf{solo sobre las primeras 76 dimensiones canónicas}, no sobre las
152. Esto es deliberado: el clipping robusto se reserva para las features continuas, no para la
máscara. El segundo es que el \textit{help text} del parser está algo desalineado con el código
ejecutado: el comentario de ayuda habla de defaults antiguos, pero \texttt{main()} resuelve hoy
\texttt{25\_000} timesteps en modo fast y \texttt{100\_000} en modo full cuando no pasas
\texttt{--timesteps}. Si te preguntan por defaults de entrenamiento, apóyate siempre en \textbf{lo
que hace el código}, no en el comentario heredado.

El mejor artefacto comprometido del repo es
\texttt{C03\_qrdqn\_cicids2017\_canonical\_full\_random\_20260223\_232439}. Su \texttt{config.json}
indica 500\,000 filas máximas cargadas, \texttt{100\,000} timesteps, \texttt{152} features,
arquitectura \texttt{[512,~256]}, \texttt{batch\_size~=~2048}, \texttt{gradient\_steps~=~20},
dispositivo \texttt{cuda} y una \texttt{reward\_config} histórica de \texttt{tp~=~1.5},
\texttt{fp~=~-2.0}, \texttt{fn~=~-5.0}, \texttt{omission~=~0.0}. Su \texttt{metrics.json} recoge
\texttt{accuracy~=~0.99859}, \texttt{recall\_attack~=~0.99945} y \texttt{f1\_attack~=~0.99876}.

Ese punto enlaza con otra respuesta de defensa imprescindible: \textbf{los defaults actuales del
código no son idénticos a la configuración del mejor run histórico}. El proyecto documenta hoy
\texttt{fp~=~-1.5} como normalización vigente, pero el mejor artefacto histórico usó
\texttt{fp~=~-2.0}. Si el tribunal te pregunta ``¿qué rewards usaba tu mejor modelo?'', no
contestes con los defaults actuales; contesta con los del artefacto concreto. Y si te preguntan
``¿qué usa hoy el código?'', entonces sí: \texttt{tp=1.5}, \texttt{fp=-1.5}, \texttt{fn=-5.0},
\texttt{omission=0.0}. Esa distinción entre \textbf{estado actual del código} y \textbf{estado
histórico del run} es uno de los mensajes más maduros metodológicamente en tu repo.

\subsection{Validación y lectura correcta de métricas}

Tu suite de validación no es un adorno; es una de las mejores partes defendibles del proyecto.
\verb|validate_checks.py| implementa tres chequeos con propósitos muy distintos. \textbf{Check~A}
evalúa el modelo de forma directa, sin depender del \texttt{info["true\_label"]} del entorno.
\textbf{Check~B} reentrena con etiquetas barajadas para detectar leakage. \textbf{Check~C} entrena y
evalúa sobre CSVs/días distintos para medir generalización dura. Además, aparte, tienes
\verb|validate_leave_one_csv_out.py|, que implementa la validación leave-one-exact-CSV-out por CSV
oficial.

Check~A sirve para responder a la objeción ``a lo mejor tu accuracy depende de cómo el entorno
informa la etiqueta''. El artefacto versionado \texttt{VAL\_checks\_A\_20260212\_235443} muestra
\texttt{accuracy~=~0.9939}, \texttt{TP~=~4772}, \texttt{FP~=~60}, \texttt{FN~=~1} y
\texttt{TN~=~5167}, usando exclusivamente predicción directa sobre \texttt{X\_test}.

Check~B es metodológicamente todavía más fuerte. El artefacto
\texttt{VAL\_checks\_B\_20260212\_235736} reporta \texttt{shuffled\_accuracy~=~0.4773}, por debajo
del baseline de clase mayoritaria \texttt{0.5227}, y deja \texttt{leakage\_detected~=~false}. La
lectura correcta es que \textbf{cuando destruyes la relación entre features y etiquetas, el
rendimiento no se mantiene artificialmente alto}. Eso es justo lo que quieres ver si no hay
leakage.

Check~C es donde el proyecto se vuelve realmente honesto. El artefacto
\texttt{VAL\_checks\_C\_20260213\_004847} baja a \texttt{accuracy~=~0.84135},
\texttt{precision\_attack~=~0.76479}, \texttt{recall\_attack~=~0.52954} y
\texttt{f1\_attack~=~0.62578}, con train en \texttt{Monday/Tuesday/Wednesday} y test en
\texttt{Thursday/Friday}. Ese salto frente al mejor resultado aleatorio (\texttt{0.99859}) es
exactamente el tipo de evidencia que demuestra que no te has quedado en un benchmark cómodo: el
pipeline funciona muy bien \textit{in-distribution}, pero generalizar a días/CSV distintos es
bastante más duro.

La validación leave-one-exact-CSV-out es, conceptualmente, aún mejor para discutir robustez porque
fuerza al modelo a no ver nunca un CSV exacto durante entrenamiento. Pero aquí la defensa honesta
importa más que el entusiasmo: el repo declara que el flujo está implementado, aunque
\textbf{todavía no hay un artefacto completo comprometido} bajo \verb|runs/validation/| para
reportar métricas consolidadas. Debes presentarlo como una validación ya soportada por el código, no
como un resultado ya cerrado.

\subsection{Fase 2, inferencia robusta y domain shift}

La Fase~2 del proyecto no es ``despliegue en producción'', sino \textbf{inferencia offline robusta
sobre tráfico capturado en un laboratorio privado}. El propio repo lo deja claro: el alcance actual
es capturar tráfico, extraer flujos, mapearlos al esquema canónico, ejecutar inferencia y guardar
artefactos reproducibles; quedan fuera el bloqueo activo en línea, la exposición pública del
laboratorio y el despliegue operacional a gran escala.

La pieza central aquí es \verb|predict_real_traffic_v2.py|, y debes saber describirla paso a paso.
El script carga un \texttt{flows.csv}, separa metadata opcional (\texttt{src\_ip},
\texttt{dst\_ip}, puertos, protocolo, timestamp y posibles columnas de verdad terreno), intenta
armonizar unidades temporales convirtiendo segundos a microsegundos cuando detecta medianas
demasiado pequeñas, mapea las columnas del extractor real a los 76 nombres canónicos, recompone el
vector de 152 dimensiones con máscara, aplica \textbf{percentile clipping} opcional sobre las
primeras 76 dimensiones, luego el \texttt{StandardScaler} persistido, luego \textbf{z-score
clipping} opcional, calcula diagnósticos de distribución, carga el modelo y predice por lotes para
evitar problemas de memoria. Finalmente guarda \texttt{predictions.csv}, \texttt{metrics.json} y
\texttt{config.json}, y opcionalmente \texttt{diagnostics.json}.

La lógica detrás de ese endurecimiento es muy defendible. \texttt{StandardScaler} es sensible a
outliers, algo que la documentación oficial de scikit-learn dice explícitamente, y el propio campo
de \textit{dataset shift} estudia precisamente qué pasa cuando la distribución de entrada en test o
despliegue difiere de la de entrenamiento. El repo introduce dos mitigaciones concretas: limitar las
features brutas a los percentiles entrenados (\(p_{0.5}\), \(p_{99.5}\)) y truncar después las
features escaladas a un máximo \(|z|\). En términos intuitivos: primero evitas que un flujo real
absurdamente extremo rompa la escala; después evitas que una z-score extrema empuje la red a un
régimen muy fuera de distribución.

Los artefactos comprometidos de Phase~2 muestran por qué esta parte todavía es científicamente
``abierta''. En \texttt{P2v2\_pred\_20260224\_004121}, sobre \verb|pcaps/flows_benign.csv|, el
sistema registró \texttt{block\_rate~=~1.0} y el \texttt{config.json} indica que \textbf{no} se
usaron percentiles (\texttt{"percentiles":~null}). En el artefacto posterior
\texttt{P2v2\_pred\_20260408\_230318}, con el mismo modelo y el mismo \texttt{scaler}, el
\texttt{block\_rate} pasó a \texttt{0.0}, y esta vez sí consta el uso de
\texttt{train\_percentiles.npz}. No debes sobrerreaccionar diciendo ``los percentiles arreglaron por
sí solos la Fase~2'', porque el repo no demuestra causalidad estricta entre ambos hechos; pero sí
puedes afirmar con rigor que la conducta de inferencia sobre tráfico benigno real \textbf{ha
cambiado sensiblemente entre artefactos}, y que una diferencia visible entre ambos runs es
precisamente el uso de clipping por percentiles.

Esa es, de hecho, la lectura metodológica correcta de la Fase~2: \textbf{el pipeline técnico está
ya bien montado; lo que sigue abierto es la robustez frente a domain shift}.

\subsection{Preguntas probables del tribunal}

\textbf{¿Por qué 76 features y no todas las del CSV original?}
Porque el repo prioriza un contrato de observación que sea \textit{flow-based}, numérico, estable y
reutilizable en tráfico real, y por eso elimina identificadores y proxies de leakage como IPs,
timestamps, Flow IDs y puertos específicos.

\textbf{¿Para qué sirve exactamente la máscara de missingness?}
Sirve para distinguir ``he imputado/esta feature no estaba disponible'' de ``he observado un valor
real''. En la implementación actual, su utilidad principal está en \textbf{mismatch estructural entre
datasets o extractores}, especialmente en Fase~2.

\textbf{¿Por qué RL y no clasificación supervisada pura?}
La respuesta honesta es que la implementación actual se parece bastante a un \textbf{contextual
bandit} o a clasificación secuencial con costes. RL sigue siendo defendible porque permite optimizar
directamente una función de utilidad asimétrica de ciberseguridad y deja el marco preparado para una
evolución futura hacia decisiones realmente secuenciales y bloqueo activo.

\textbf{¿Qué optimiza de verdad el sistema de recompensas?}
Optimiza una asimetría de seguridad: castiga mucho más dejar pasar ataques (\texttt{fn~=~-5.0}) que
bloquear benignos (\texttt{fp~=~-1.5}), recompensa bloquear ataques (\texttt{tp~=~1.5}) y deja
neutro permitir benignos (\texttt{omission~=~0.0}).

\textbf{¿Por qué no basta con la accuracy del mejor modelo?}
Porque el mejor resultado comprometido (\texttt{0.99859}) viene de un split aleatorio sobre la misma
distribución; cuando pasas a una validación por días/CSV distintos, la accuracy baja a
\texttt{0.84135} y el recall de ataque a \texttt{0.5295}.

\textbf{¿Cómo has intentado evitar leakage?}
De tres formas: quitando features potencialmente espurias; ajustando el \texttt{scaler} solo sobre
train; y ejecutando un shuffled-label test donde el rendimiento no se mantiene artificialmente alto.

\textbf{¿Qué parte está realmente resuelta y cuál no?}
Está resuelta la arquitectura experimental: contrato de datos, adapter, entorno, entrenamiento
reproducible, validación A/B/C y pipeline robusto de Fase~2. No está resuelto todavía el
\textbf{bloqueo activo en tiempo real} ni la \textbf{robustez cerrada frente a tráfico real benigno
y mixto}.

\textbf{¿Cuál es la formulación más sólida para cerrar la defensa?}
Que tu TFG no aporta solo ``un modelo que clasifica'', sino \textbf{un pipeline de defensa RL
reproducible y trazable}, con un contrato de observación coherente, control explícito del leakage,
validación dura y una transición ya operativa hacia tráfico real offline.

\end{document}
